\section*{Exercises}
\begin{exercise}\label{ex:12.1}Let $T(x)=(3x+1)/4$ on $\RR$. Find the fixed point of $T$ and the smallest contraction factor of $T$. Then find a formula for the iterates $x_n$ starting from $x_0=0$, and prove it by induction.\end{exercise}
\begin{exercise}\label{ex:12.2}Let $T$ be a contraction with factor $q=3/4$ of a nonempty complete metric space, with fixed point $x_*$, and let $z$ be a point with $d(z,Tz)\le0.002$. What upper bound on $d(z,x_*)$ does Theorem~\ref{thm:banach} guarantee? Is the residual bound $0.002$ itself an upper bound on $d(z,x_*)$ in every such case? Give a reason or an example.\end{exercise}
\begin{exercise}\label{ex:12.3}Let $T(x)=(x^2+1)/3$. Show that $T$ maps the interval $[1/3,2/3]$ into itself and that $4/9$ is a contraction factor of $T$ on this interval. Starting from $x_0=0$, explain why every iterate from $x_1$ on lies in $[1/3,2/3]$. Use the smaller factor to improve the certificate $d(x_2,x_*)\le19/729$ obtained in Section~\ref{ch12:sec:example}, and explain why the improvement is justified.\end{exercise}
\begin{exercise}\label{ex:12.4}Section~\ref{ch12:sec:hypotheses} gives counterexamples in three groups: without completeness, without a factor below one, and without the self-map condition. For each example in these three groups, state which hypothesis of Theorem~\ref{thm:banach} fails, and prove that all the other hypotheses hold. For $T(x)=x+1/x$, this includes proving that $[1,\infty)$ is complete.\end{exercise}
\begin{exercise}\label{ex:12.5}Let $a$ and $b$ be real numbers with $|a|<1$, and let $T(x)=ax+b$ on $\RR$. Find the fixed point $x_*$ of $T$, and prove by induction that for every starting point $x_0$ the iterates satisfy $x_n-x_*=a^n(x_0-x_*)$ for every $n\in\NN$. How do the iterates move when $a<0$, compared with when $0<a<1$?\end{exercise}
\begin{exercise}\label{ex:12.6}Let $T$ and $S$ be two self-maps of the same nonempty complete metric space $(X,d)$, both contractions with factor $q$, where $0\le q<1$. Write $x_T$ and $x_S$ for their fixed points. Suppose there is a number $\eta\ge0$ with $d(Tx,Sx)\le\eta$ for every $x\in X$. Prove that $d(x_T,x_S)\le\eta/(1-q)$.\end{exercise}
